%% kapa: exam=1 round=32 year=2021 period=1 form=A subject=civil_law subject_ko=민법 seq=22 number=22
%% source: Q-Net 감정평가사 기출문제 (cst00309) artlSeq=5208397 file=extracted/1008/5208397/A형/2021년도 제32회 감정평가사 1차 1교시 A형.pdf
%% conversion: kordoc 4.24.51 -> tools/gen1.py (auto); review=proofread-sonnet; flags=none
\begin{question}{22}
점유자와 회복자의 관계에 관한 설명으로 옳은 것은? (다툼이 있으면 판례에 따름)
\begin{choices}
  \item 선의의 점유자가 취득하는 과실에 점유물의 사용이익은 포함되지 않는다.
  \item 유치권자에게는 원칙적으로 수익목적의 과실수취권이 인정된다.
  \item 점유물이 점유자의 귀책사유로 훼손된 경우, 선의의 점유자는 소유의 의사가 없더라도 이익이 현존하는 한도에서 배상책임이 있다.
  \item 회복자로부터 점유물의 반환을 청구 받은 점유자는 유익비의 상환을 청구할 수 있다.
  \item 점유물의 소유자가 변경된 경우, 점유자는 유익비 지출 당시의 전 소유자에게 비용의 상환을 청구해야 한다.
\end{choices}
\end{question}
\answer{4}
